\Needspace{10\baselineskip}
\section{Indexing}

\subsubsection{ব্যাংকের তথ্য অনুসন্ধানে Indexing}
ব্যাংকে AccountNo দিয়ে বারবার নির্দিষ্ট হিসাব খুঁজতে হয়। সহায়ক ব্যবস্থা না থাকলে প্রতিবার অনেক রেকর্ড পরীক্ষা করতে হতে পারে। 

ভর্তি ব্যবস্থায় RegistrationNo, লাইব্রেরিতে BookID এবং অনলাইন অর্ডারে OrderID দিয়ে নিয়মিত খোঁজার ক্ষেত্রেও ইনডেক্স কাজে লাগতে পারে।

\begin{definitionbox}{Indexing (ইনডেক্সিং)}
এক বা একাধিক ফিল্ডের মান থেকে সংশ্লিষ্ট রেকর্ড দ্রুত খুঁজতে সহায়ক নির্দেশনামূলক কাঠামো তৈরি ও রক্ষণাবেক্ষণের প্রক্রিয়াকে ইনডেক্সিং বলে।
\end{definitionbox}
বইয়ের নির্ঘণ্টে বিষয় দেখে পৃষ্ঠা খুঁজে পাওয়া যায়। ডেটাবেজের \textbf{Index}-ও ফিল্ডের মানের সঙ্গে রেকর্ডের নির্দেশনা রাখে। এতে অনেক ক্ষেত্রে প্রতিবার সব রেকর্ড পরীক্ষা করতে হয় না। ইনডেক্স মূল হিসাবের বিকল্প তালিকা নয়। 

\subsection{Index-এর কাজ ও নির্বাচন}
হাজার শিক্ষার্থীর মধ্যে Roll 103 খুঁজতে Roll-ভিত্তিক ইনডেক্স সহায়তা করতে পারে। Name, RegistrationNo বা AccountNo দিয়ে বারবার অনুসন্ধান করা হলে সংশ্লিষ্ট ফিল্ডেও ইনডেক্স বিবেচনা করা যায়। তবে কুয়েরি পরিকল্পনাকারী প্রয়োজন অনুযায়ী ইনডেক্স ব্যবহার করে; ইনডেক্স থাকলেই প্রতিটি কুয়েরি দ্রুত হবে না।
\begin{flushleft}\small
\begin{tabularx}{\linewidth}{|p{30mm}|X|}\hline
বিবেচনা & ইনডেক্স নির্বাচনের যুক্তি\\\hline
অনুসন্ধানের ব্যবহার & যে ফিল্ডে বারবার শর্ত বা সংযোগ দেওয়া হয় তা বিবেচ্য\\\hline
মানের বৈচিত্র্য & অনেক পৃথক মান থাকলে নির্দিষ্ট রেকর্ড খোঁজা ফলপ্রসূ হতে পারে\\\hline
রেকর্ড পরিবর্তন & ঘন ঘন লেখা হলে ইনডেক্স রক্ষণাবেক্ষণের ব্যয় বাড়ে\\\hline
সংরক্ষণের জায়গা & প্রতিটি ইনডেক্সের জন্য অতিরিক্ত স্থান প্রয়োজন\\\hline
\end{tabularx}
\end{flushleft}
ইনডেক্স তৈরি করলে মূল ফিল্ডের অর্থ বা শিক্ষার্থীর সংরক্ষিত GPA বদলায় না। 


\subsubsection{Index-এর বৈশিষ্ট্য}
\begin{enumerate}
\item \textbf{সহায়ক কাঠামো:} ফিল্ডের মানের সাহায্যে সংশ্লিষ্ট রেকর্ড খুঁজতে সহায়তা করে।
\item \textbf{হালনাগাদ:} রেকর্ড যোগ, পরিবর্তন বা মুছে ফেললে সংশ্লিষ্ট Index-ও হালনাগাদ করতে হয়।
\end{enumerate}
\textbf{প্রয়োগ:} SELECT ছাড়াও UPDATE, DELETE ও JOIN-এর জন্য প্রয়োজনীয় রেকর্ড খুঁজতে Index কাজে লাগতে পারে। কতটা সুবিধা হবে, তা শর্তে মেলা রেকর্ডের সংখ্যা ও Index-এর গঠনের ওপর নির্ভর করে। অল্প রেকর্ড প্রয়োজন হলে Index সুবিধাজনক হতে পারে; অধিকাংশ রেকর্ড লাগলে DBMS অন্য পদ্ধতি বেছে নিতে পারে।
\subsubsection{সুবিধা}
\begin{enumerate}
\item \textbf{দ্রুত অনুসন্ধান:} নিয়মিত Roll দিয়ে নির্দিষ্ট শিক্ষার্থীর ফি খোঁজার মতো কাজে সহায়তা করে।
\end{enumerate}
\subsubsection{অসুবিধা}
\begin{enumerate}
\item \textbf{অতিরিক্ত স্থান:} Index সংরক্ষণে বাড়তি জায়গা লাগে।
\item \textbf{রক্ষণাবেক্ষণের ব্যয়:} ডেটা পরিবর্তনের সঙ্গে Index পরিবর্তনের কাজও করতে হয়।
\end{enumerate}
\textbf{মনে রাখবে:} একবার পুরো তালিকা পড়ার জন্য বহু Index বানালে লাভের তুলনায় ব্যয় বেশি হতে পারে। প্রয়োজনীয় অনুসন্ধান ও পরিবর্তনের কাজ মিলিয়ে Index নির্বাচন করতে হয়।

\subsection{Index-এর ধরন}
\textbf{Automatic ও Manual:} কোনো কী বা বিধি নির্ধারণের ফলে DBMS নিজে সহায়ক ইনডেক্স তৈরি করলে তাকে স্বয়ংক্রিয় ইনডেক্সের উদাহরণ বলা যায়। ব্যবহারকারী প্রয়োজনমতো নির্দেশ বা নকশা-পরিবেশ থেকে তৈরি করলে তা ম্যানুয়াল ইনডেক্স। কোন বিধিতে কী ইনডেক্স স্বয়ংক্রিয় হয়, তা DBMS-নির্ভর।

\textbf{Unique ও Non-unique:} Unique index একই কী-মানের পুনরাবৃত্তি প্রতিরোধ করে। RegistrationNo অনন্য রাখার জন্য এটি উপযোগী। Name পুনরাবৃত্ত হতে পারলে non-unique index ব্যবহার করা যায়। Unique index-এ NULL-এর আচরণ সব DBMS-এ এক নয়; তাই unique বলেই সেটিকে প্রাইমারি কী-এর সব নিয়মের সমান ধরা যাবে না।

\textbf{Single-field ও Composite:} একটি ফিল্ডের ওপর তৈরি ইনডেক্সকে Single-field Index এবং একাধিক ফিল্ডের ওপর তৈরি ইনডেক্সকে Composite বা Compound Index বলে। ClassName ও Roll-এর যৌগিক unique index থাকলে একই শ্রেণি ও রোলের জোড়া দ্বিতীয়বার গ্রহণযোগ্য নয়; ভিন্ন শ্রেণিতে একই রোল থাকতে পারে। যদি একই শ্রেণি বিভিন্ন শিক্ষাবর্ষে পুনরায় থাকে, তখন অনন্যতার নিয়মে Session-ও বিবেচ্য।

\subsection{Access-এর Indexed Property}
Table Design View-তে ফিল্ড নির্বাচন করে Indexed property নির্ধারণ করা যায়। \textbf{No} মানে ঐ বৈশিষ্ট্য দিয়ে ইনডেক্স নির্ধারিত নয়; \textbf{Yes (Duplicates OK)} একই মান গ্রহণ করে এবং \textbf{Yes (No Duplicates)} একই মানের পুনরাবৃত্তি প্রতিরোধ করে। Name-এর জন্য প্রথম Yes এবং অনন্য পরিচিতির জন্য দ্বিতীয় Yes যৌক্তিক হতে পারে। একাধিক ফিল্ডের ইনডেক্স তৈরিতে Indexes উইন্ডোতে একই ইনডেক্সের অধীনে ফিল্ডগুলো দেওয়া যায়।

% ch06_access_indexed_property.png is not used: it shows Indexed=No
% for a field marked as a primary key. Keep the accurate explanation above.

\subsection{SQL দিয়ে Index তৈরি ও অপসারণ}
নিচের SQLite উদাহরণে Student টেবিল আগে থেকেই আছে এবং তাতে Name, RegistrationNo, ClassName ও Roll ফিল্ড আছে ধরা হয়েছে। প্রতিটি নির্দেশ আলাদা উদ্দেশ্যে ব্যবহৃত।
\begin{verbatim}
CREATE INDEX idx_student_name ON Student(Name);
CREATE UNIQUE INDEX idx_student_reg
ON Student(RegistrationNo);
CREATE UNIQUE INDEX idx_class_roll
ON Student(ClassName, Roll);
DROP INDEX idx_student_name;
\end{verbatim}
শেষ নির্দেশটি শুধু ইনডেক্স বাদ দেয়, Student টেবিল বা শিক্ষার্থীর রেকর্ড নয়। unique index তৈরির আগে বিদ্যমান ডেটায় একই মান থাকলে নির্দেশ ব্যর্থ হতে পারে; আগে অনন্যতার নিয়ম ও ডেটা মিলিয়ে নিতে হবে। অন্য DBMS-এ DROP INDEX-এর গঠন ভিন্ন হতে পারে।

\subsection{Sorting, Indexing ও Key-এর পার্থক্য}
\begin{flushleft}\small
\begin{tabularx}{\linewidth}{|p{26mm}|X|X|}\hline
বিষয় & সর্টিং & ইনডেক্সিং\\\hline
প্রধান উদ্দেশ্য & ফল নির্দিষ্ট ক্রমে সাজানো & অনুসন্ধানের সহায়ক কাঠামো তৈরি\\\hline
SQL উদাহরণ & ORDER BY & CREATE INDEX\\\hline
অতিরিক্ত কাঠামো & স্থায়ী ইনডেক্স আবশ্যক নয় & ইনডেক্সের স্থান ও রক্ষণাবেক্ষণ লাগে\\\hline
ফলক্রম & নির্দেশ অনুযায়ী নির্ধারিত & নিজে ফলের ক্রমের নিশ্চয়তা নয়\\\hline
\end{tabularx}
\end{flushleft}
প্রাইমারি কী রেকর্ডের অনন্য পরিচয় নির্ধারণ করে এবং শুদ্ধতার নিয়ম বজায় রাখে। ইনডেক্স অনুসন্ধানে সহায়তার জন্য তৈরি কাঠামো। কী বাস্তবায়নে ইনডেক্স ব্যবহৃত হতে পারে, কিন্তু সব ইনডেক্স প্রাইমারি কী নয়। একই কাজে কুয়েরি প্রয়োজনীয় ডেটা বেছে নেয়, সর্টিং ফল সাজায় এবং ইনডেক্স অনুসন্ধানে সহায়তা করতে পারে।

\subsection{মূল Record ও Index-এর উদাহরণ}
নিচের সারণিতে R1, R2 ইত্যাদি দিয়ে রেকর্ডের অবস্থান বোঝানো হয়েছে; এটি কোনো নির্দিষ্ট DBMS-এর প্রকৃত ডিস্ক ঠিকানার ছক নয়। মূল রেকর্ডের বর্তমান অবস্থান এবং Roll-এর সংগঠিত ইনডেক্স পাশাপাশি দেখো।
\begin{flushleft}\small
\begin{tabular}{|l|l|l|}\hline
নির্দেশনা & Roll & Name\\\hline
R1 & 104 & Hasan\\\hline
R2 & 101 & Rupa\\\hline
R3 & 103 & Tania\\\hline
R4 & 102 & Nabil\\\hline
\end{tabular}\qquad
\begin{tabular}{|l|l|}\hline
Roll & রেকর্ডের নির্দেশনা\\\hline
101 & R2\\\hline102 & R4\\\hline103 & R3\\\hline104 & R1\\\hline
\end{tabular}
\end{flushleft}
\begin{flushleft}
\begin{tikzpicture}[every node/.style={font=\small}]
\node[booknode,text width=30mm] (q) {সন্ধান: Roll 103};
\node[booknode,text width=31mm,right=7mm of q] (i) {ইনডেক্স: 103 $\to$ R3};
\node[booknode,text width=31mm,right=7mm of i] (r) {R3: 103, Tania};
\draw[bookarrow] (q)--(i);\draw[bookarrow] (i)--(r);
\end{tikzpicture}\end{flushleft}
\begin{boardbox}{ইনডেক্স ও ফলের ক্রম}
ইনডেক্স সন্ধানে সহায়তা করতে পারে; ফলের নির্দিষ্ট ক্রম চাইলে ORDER BY দিতে হয়। অতিরিক্ত ইনডেক্সে জায়গা লাগে এবং ডেটা বদলালে ইনডেক্সও হালনাগাদ করতে হয়।
\end{boardbox}
Roll 103 খুঁজলে ইনডেক্সের মাধ্যমে R3-এর রেকর্ডে পৌঁছানো যায়। মূল রেকর্ডের Name বা Roll পরিবর্তিত হয়নি। যদি Roll 105-এর নতুন রেকর্ড যোগ হয়, তবে উপযুক্ত ইনডেক্স কাঠামোতেও তার মান ও নির্দেশনা যুক্ত করতে হয়। ইনডেক্স তৈরি করার পর ব্যবহারকারীকে প্রতিবার হাতে এই ছক লেখা লাগে না; DBMS রক্ষণাবেক্ষণ করে।

\subsection{Index ফাইলের পুরোনো পরিভাষা}
পুরোনো FoxPro-ধরনের ব্যবস্থায় স্বতন্ত্র single index-এর জন্য \texttt{.idx} এবং compound index file-এর জন্য \texttt{.cdx} নাম দেখা যায়। একটি compound file-এ একাধিক নামযুক্ত index tag থাকতে পারে। ফাইলের compound পরিচয় এবং এক ইনডেক্সের composite key একেবারে একই বিষয় নয়।

\texttt{.idx} ও \texttt{.cdx} নির্দিষ্ট পুরোনো সফটওয়্যারের ফাইল এক্সটেনশন। Access বা SQLite-এর প্রতিটি ইনডেক্স এমন আলাদা ফাইলে থাকে না। ইনডেক্সিংয়ের সাধারণ বৈশিষ্ট্য হলো অনুসন্ধানে সহায়তা করা; ওই ফাইল এক্সটেনশন সব DBMS-এর বৈশিষ্ট্য নয়।

\subsection{একই ফিল্ডে Index ও Key-এর পৃথক নিয়ম}
Student-এর Name-এ Rupa দুইবার এবং RegistrationNo-তে R501 ও R502 আছে। নাম দিয়ে খুঁজতে non-unique index উপযোগী হতে পারে। RegistrationNo প্রতিষ্ঠানে অনন্য হলে unique index প্রাসঙ্গিক। Name-এ No Duplicates দিলে দ্বিতীয় Rupa-র এন্ট্রি ভুলভাবে আটকে যাবে। আবার RegistrationNo-তে Duplicates OK দিলে প্রত্যাশিত অনন্যতার নিয়মটি ঐ ইনডেক্স দিয়ে প্রয়োগ হবে না।

একটি ইনডেক্স দ্রুত পড়ায় সহায়তা করলেও রেকর্ড লেখা ও বদলানোর ব্যয় বাড়ায়। প্রতিদিন লক্ষবার অনুসন্ধান ও অল্প পরিবর্তনের ব্যবস্থায় একটি ইনডেক্স উপকারী হতে পারে; ছোট টেবিলে অব্যবহৃত বহু ইনডেক্স থেকে সমপরিমাণ লাভ নাও পাওয়া যেতে পারে। প্রয়োজনীয় কাজ ও ডেটার প্রকৃতি মিলিয়ে সিদ্ধান্ত দিতে হয়।

\subsection{পণ্যের নামে অনুসন্ধান ও মূল্যে সাজানো}
\textbf{নমুনা।} একটি সমবায়ে Product(ProductID, ProductName, UnitPrice, Stock) টেবিল আছে। ProductID অনন্য। কর্মীরা পণ্যের নাম দিয়ে খোঁজেন। ব্যবস্থাপক বেশি দাম থেকে কম দামে তালিকা চান। বিক্রয়কর্মী মজুদ দেখতে পারবেন, কিন্তু মূল্য পরিবর্তনের অনুমতি পাবেন না।

\textbf{ইনডেক্স ও সাজানো ফলের নির্দেশ:}
\begin{verbatim}
CREATE INDEX idx_product_name ON Product(ProductName);
SELECT ProductID, ProductName, UnitPrice
FROM Product
ORDER BY UnitPrice DESC, ProductID ASC;
\end{verbatim}
প্রথম নির্দেশে ইনডেক্স তৈরি হয়; রেকর্ডের সংখ্যা বদলায় না। দ্বিতীয়টি নির্বাচিত কলামের তালিকা দেখায়; সমান মূল্যে ProductID দিয়ে ক্রম ঠিক হয়। পণ্যের নমুনা মান না দেওয়ায় নির্দিষ্ট নামের আউটপুট বানানো যাবে না।

\textbf{একই ব্যবস্থায় অধিকার নিয়ন্ত্রণ:} বিক্রয়কর্মীর মূল্য পরিবর্তনের অনুমতি বন্ধ রেখে প্রয়োজনমতো দেখা বা বিক্রয় এন্ট্রির সুযোগ দিতে হবে। অনুমোদিত ব্যবস্থাপক মূল্য বদলালে কে কখন বদলালেন তার লগ রাখা হবে। ভুল পরিবর্তন বা ডেটা হারালে ব্যাকআপ ও পুনরুদ্ধার দরকার। শুধু পাসওয়ার্ড থাকলেই এই তিন কাজ একসঙ্গে সম্পন্ন হয় না।

\subsection{গুরুত্বপূর্ণ প্রশ্ন}
\textbf{অনুশীলন:} (১) একই নামের বহু ব্যক্তি থাকলে Name-এ কোন ধরনের ইনডেক্স উপযুক্ত? (২) ClassName ও Roll-এর unique index থাকলে XI, 1 ও XII, 1 গ্রহণযোগ্য কি? (৩) একটি ইনডেক্স বাদ দিলে কেন টেবিলের রেকর্ড থাকে?\par



